
\newcommand{\mectable}[4]{%
  \begin{table}[H]
    \centering
    \caption{#1}
    \label{#2}
    \resizebox{\textwidth}{!}{\input{#3}}
    {\footnotesize\raggedright\noindent Notes: #4\par}
  \end{table}
  \clearpage
}

\begin{table}[H]
\RawCaption%
  {\caption{\deleted{Barham and Rowberry’s results and our replication: Impact of Progresa on elderly mortality 65+, municipalities incorporated in 1998 or 1999}\added{Replication of Barham and Rowberry (2013): Effect of Progresa on Elder Mortality Rate (EMR 65+), Municipalities Incorporated in 1998--1999}}
   \label{t:did_robust_br}}
\floatbox[{\captop}]{table}[\FBwidth]
  {}
  {\scalebox{0.9}{
      \begin{tabular}{lccc} \hline \hline
& \multicolumn{1}{c}{Pooled} & \multicolumn{1}{c}{Males} & \multicolumn{1}{c}{Females} \\ 
\cmidrule(lr){2-2}\cmidrule(lr){3-3}\cmidrule(lr){4-4}  
& \multicolumn{1}{c}{(1)} & \multicolumn{1}{c}{(2)} & \multicolumn{1}{c}{(3)} \\ \toprule 
\underline{\textit{Panel A: BR (2013)}}  \\ 
\textit{2-yr lagged Intensity} & -6.37*** & -6.42*** & -6.46*** \\ 
 & (1.04) & (1.42) & (1.31) \\ 
  & & & \\ 
Mean 1996 & 47.5 & 49.3 & 46.0 \\ 
Obs & 21,571 & 21,571 & 21,571 \\ 
No. Mun & 1,961 & 1,961 & 1,961 \\ 
  & & & \\ 
\underline{\textit{Panel B: Replication (Unweighted)}}  \\ 
\textit{2-yr lagged Intensity} &       -7.061*** &       -6.889*** &       -7.426*** \\ 
 & (       1.256) & (       1.711) & (       1.559) \\ 
  & & & \\ 
\underline{\textit{Panel C: Replication (Weighted)}}  \\ 
\textit{2-yr lagged Intensity} &       -3.740*** &       -3.322*** &       -4.210*** \\ 
 & (       0.786) & (       0.935) & (       1.002) \\ 
  & & & \\ 
\underline{\textit{Panel D: 1-yr Lag (Weighted)}}  \\ 
\textit{1-yr lagged Intensity} &       -3.457*** &       -2.681** &       -4.255*** \\ 
 & (       0.890) & (       1.072) & (       1.043) \\ 
  & & & \\ 
\underline{\textit{Panel E: 3-yr Lag (Weighted)}}  \\ 
\textit{3-yr lagged Intensity} &       -2.071*** &       -2.357** &       -1.830* \\ 
 & (       0.772) & (       0.957) & (       0.960) \\ 
  & & & \\ 
Mean 1996 &        48.16 &        48.87 &        47.49 \\ 
Obs &       15,642 &       15,642 &       15,642 \\ 
No. Mun &        1,422 &        1,422 &        1,422 \\ 
  & & & \\ 
Year FE & Y & Y & Y \\ 
Mun FE & Y & Y & Y \\ 
Cluster SE: Mun & Y & Y & Y \\ 
\bottomrule
\end{tabular}
    }
    \vspace{-3mm}
    \floatfoot{\deleted{Notes: The data analysis is restricted to municipalities incorporated into Progresa in 1998 or 1999. Progresa intensity data are pooled from 1990 to 2000. All regression models include year fixed effects and municipality fixed effects. Standard errors are clustered at the municipality level. The regression model is weighted by the number of population aged +65 for both sexes in (1), only male in (2), and only female in (3).\\ ***p$<0.01$, **p$<0.05$, *p$<0.1$.}\added{Notes: This table replicates the main specification of Barham and Rowberry (2013). The sample is restricted to municipalities incorporated into Progresa in 1998 or 1999. The outcome is the Elder Mortality Rate (EMR) for adults aged 65 and older per 1,000 inhabitants. Progresa intensity is measured using data pooled from 1990 to 2000. All regressions include municipality and year fixed effects. The regression model is weighted by the population aged 65 and older, shown separately for the pooled (column 1), male (column 2), and female (column 3) populations. Standard errors are clustered at the municipality level and shown in parentheses. ***p$<0.01$, **p$<0.05$, *p$<0.1$.}}
  }
\end{table}
\newpage




\begin{table}[H]
\RawCaption%
  {\caption{\deleted{Barham and Rowberry Extended Results Replication}\added{Sensitivity Analysis: Barham and Rowberry (2013) Replication with Alternative Lag Structures and Sample Restrictions}}
   \label{t:did_robust_br2}}
\floatbox[{\captop}]{table}[\FBwidth]
  {}
  {\scalebox{0.9}{
      \input{tables/appendix/T_BR_robustness_emr65}
    }
    \vspace{-3mm}
    \floatfoot{\deleted{Notes: Progresa intensity data are pooled from 1990 to 2000. All regression models include year fixed effects and municipality fixed effects. Standard errors are clustered at the municipality level. The regression model is weighted by the number of population aged +65. ***p$<0.01$, **p$<0.05$, *p$<0.1$.}\added{Notes: This table reports sensitivity analyses extending the replication of Barham and Rowberry (2013), varying the lag structure of Progresa intensity, weighting, and sample restrictions. The outcome is the Elder Mortality Rate (EMR) for adults aged 65 and older per 1,000 inhabitants. Row (A) reproduces Barham and Rowberry's (2013) original results; row (B) reports our replication using the same specification. Subsequent rows vary the lag structure (1-year and 3-year lags), weighting, and sample. Progresa intensity data are pooled from 1990 to 2000. All regressions include municipality and year fixed effects. Standard errors are clustered at the municipality level. ***p$<0.01$, **p$<0.05$, *p$<0.1$.}}
  }
\end{table}

% \begin{table}[H]
% \RawCaption%
%   {\caption{Barham and Rowberry Extended Results Replication}
%    \label{t:did_robust}}
% \floatbox[{\captop}]{table}[\FBwidth]
%   {}
%   {\scalebox{0.9}{
%       \input{tables/appendix/T_BR_robustness_aamr65}
%     }
%     \vspace{-3mm}
%     \floatfoot{Notes: Progresa intensity data are pooled from 1990 to 2000. All regression models include year fixed effects and municipality fixed effects. Standard errors are clustered at the municipality level. The regression model is weighted by the number of population aged +65. ***p$<0.01$, **p$<0.05$, *p$<0.1$.}
%   }
% \end{table}





% \begin{table}[H]
% \RawCaption%
%   {\caption{Functional Forms}
%    \label{t:did_robust}}
% \floatbox[{\captop}]{table}[\FBwidth]
%   {}
%   {\scalebox{0.9}{
%       \begin{tabular}{lccc} \hline \hline
& Levels & Log & Poisson \\ 
\cmidrule(lr){2-2}\cmidrule(lr){3-3}\cmidrule(lr){4-4}
& \multicolumn{1}{c}{(1)} & \multicolumn{1}{c}{(2)} & \multicolumn{1}{c}{(3)} \\ \toprule
\underline{\textit{Panel A: Pooled}}  \\ 
\textit{Intensity 1999 x post (1997-2006)} &       -0.777 &       -0.020 &       -0.031 \\ 
  & (       1.464) & (       0.042) & (       0.058) \\ 
   & & & \\ 
\textit{Intensity 2005 x post (1997-2006)} &       -2.787 &       -0.052 &       -0.128* \\ 
 & (       1.806) & (       0.052) & (       0.072) \\ 
  & & & \\ 
Mean (1991-1996)  &        46.47 &         3.75 &        23.86 \\ 
  & & &  \\ 
\underline{\textit{Panel B: Males}}  \\ 
\textit{Intensity 1999 x post (1997-2006)} &       -2.288 &       -0.029 &       -0.104* \\ 
  & (       1.732) & (       0.050) & (       0.062) \\ 
  & & & \\ 
\textit{Intensity 2005 x post (1997-2006)}  &        0.214 &       -0.020 &       -0.024 \\ 
 & (       2.000) & (       0.057) & (       0.076) \\ 
 & & & \\ 
Mean (1991-1996) &        47.26 &         3.78 &        12.25 \\ 
  & & &  \\ 
\underline{\textit{Panel C: Females}}  \\ 
\textit{Intensity 1999 x post (1997-2006)}  &        0.592 &        0.025 &        0.055 \\ 
  & (       1.824) & (       0.055) & (       0.074) \\ 
  & & & \\ 
\textit{Intensity 2005 x post (1997-2006)} &       -5.637** &       -0.173** &       -0.223*** \\ 
 & (       2.205) & (       0.067) & (       0.076) \\ 
   & & & \\ 
Mean (1991-1996)  &        46.11 &         3.73 &        11.61 \\ 
  & & &  \\ 
Obs &       17,904 &       17,223 &       17,904 \\ 
No. Mun &        1,119 &        1,119 &        1,119 \\ 
  & & & \\ 
Year FE & Y & Y & Y \\ 
Mun FE  & Y & Y & Y \\ 
Weights & Y & Y & Y \\ 
Cluster SE: Mun & Y & Y & Y \\ 
\bottomrule
\end{tabular}
%     }
%     \vspace{-3mm}
%     \floatfoot{Notes: Progresa intensity data are pooled from 1990 to 2000. All regression models include year fixed effects and municipality fixed effects. Standard errors are clustered at the municipality level. The regression model is weighted by the number of population aged +65. ***p$<0.01$, **p$<0.05$, *p$<0.1$.}
%   }
% \end{table}
% \newpage






\begin{table}[H]
\RawCaption%
  {\caption{\deleted{AAMR}\added{Effects of Progresa Enrollment Intensity on Age-Adjusted Mortality Rate (AAMR 65+) by Sex, Marginalized Municipalities, 1991--2006}}
   \label{t:did_robust_aamr}}
\floatbox[{\captop}]{table}[\FBwidth]
  {}
  {\scalebox{0.9}{
      \input{tables/appendix/AT4_aamr_mortality}
    }
    \vspace{-3mm}
    \floatfoot{\deleted{Notes: Progresa intensity data are pooled from 1990 to 2000. All regression models include year fixed effects and municipality fixed effects. Standard errors are clustered at the municipality level. The regression model is weighted by the number of population aged +65. ***p$<0.01$, **p$<0.05$, *p$<0.1$.}\added{Notes: This table reports difference-in-differences estimates of the effect of Progresa enrollment intensity on the Age-Adjusted Mortality Rate (AAMR) for adults aged 65 and older per 1,000 inhabitants, separately for the pooled, male, and female populations. Columns (1) and (2) report unweighted estimates without and with Seguro Popular intensity; columns (3) and (4) report estimates weighted by the population aged 65 and older without and with Seguro Popular intensity. The post-dummy equals 1 for years 1997--2006. The sample is restricted to highly marginalized municipalities ($\textit{gm\_mun\_1990} = 4$ or $5$). All regressions include municipality and year fixed effects. Standard errors are clustered at the municipality level. ***p$<0.01$, **p$<0.05$, *p$<0.1$.}}
  }
\end{table}

\begin{landscape}
\begin{table}[H]
\RawCaption%
  {\caption{\deleted{AAMR}\added{Effects of Progresa Enrollment Intensity on Cause-Specific Elder Mortality Rate (EMR 65+), Highly Marginalized Municipalities, 1991--2006}}
   \label{t:did_robust_cod}}
\floatbox[{\captop}]{table}[\FBwidth]
  {}
  {\scalebox{0.7}{
      \input{tables/appendix/AT5_cod_mortality}
    }
    \vspace{-3mm}
    \floatfoot{\deleted{Notes: Progresa intensity data are pooled from 1990 to 2000. All regression models include year fixed effects and municipality fixed effects. Standard errors are clustered at the municipality level. The regression model is weighted by the number of population aged +65. ***p$<0.01$, **p$<0.05$, *p$<0.1$.}\added{Notes: This table reports difference-in-differences estimates of the effect of Progresa enrollment intensity on cause-specific Elder Mortality Rate (EMR) for adults aged 65 and older per 1,000 inhabitants. Each column corresponds to a specific cause of death. The post-dummy equals 1 for years 1997--2006. The sample is restricted to highly marginalized municipalities ($\textit{gm\_mun\_1990} = 4$ or $5$). All regressions include municipality and year fixed effects. Standard errors are clustered at the municipality level. ***p$<0.01$, **p$<0.05$, *p$<0.1$.}}
  }
\end{table}
\end{landscape}

\begin{landscape}
\begin{table}[H]
\RawCaption%
  {\caption{\deleted{AAMR}\added{Effects of Progresa Enrollment Intensity on Cause-Specific Elder Mortality Rate (EMR 65+), Barham-Rowberry Sample, 1992--2002}}
   \label{t:did_robust_cod_br}}
\floatbox[{\captop}]{table}[\FBwidth]
  {}
  {\scalebox{0.9}{
      \begin{tabular}{lcccccccccc} \hline \hline
& Cancer & Diab. & IllDef & Resp. & Card. & Infect. & Nutri. & Accid. & Other \\ 
\cmidrule(lr){2-2}\cmidrule(lr){3-3}\cmidrule(lr){4-4}\cmidrule(lr){5-5}\cmidrule(lr){6-6}\cmidrule(lr){7-7}\cmidrule(lr){8-8}\cmidrule(lr){9-9}\cmidrule(lr){10-10}
& \multicolumn{1}{c}{(1)} & \multicolumn{1}{c}{(2)} & \multicolumn{1}{c}{(3)} & \multicolumn{1}{c}{(4)} & \multicolumn{1}{c}{(5)} & \multicolumn{1}{c}{(6)} & \multicolumn{1}{c}{(7)} & \multicolumn{1}{c}{(8)} & \multicolumn{1}{c}{(9)} \\ \toprule
\underline{\textit{Panel A: Pooled}}  \\ 
\textit{Intensity 1999 x post} &       -0.816*** &       -0.924*** &       -0.614*** &        0.149 &        0.792* &       -1.670*** &       -0.279 &        0.019 &        1.681*** \\ 
 & (       0.188) & (       0.183) & (       0.134) & (       0.281) & (       0.470) & (       0.210) & (       0.248) & (       0.048) & (       0.498) \\ 
\textit{RW p-value} & [ 0.000] & [ 0.000] & [ 0.000] & [ 1.000] & [ 0.370] & [ 0.000] & [ 0.784] & [ 1.000] & [ 0.004] \\ 
  & & & & & & & & & \\ 
Mean (pre-1997) &         2.48 &         2.58 &         0.91 &         4.86 &        15.84 &         2.56 &         3.54 &         0.27 &        14.39 \\ 
Obs &       15,642 &       15,642 &       15,642 &       15,642 &       15,642 &       15,642 &       15,642 &       15,642 &       15,642 \\ 
  & & & & & & & & & \\ 
\underline{\textit{Panel B: Females}}  \\ 
\textit{Intensity 1999 x post} &       -0.600*** &       -0.758*** &       -0.687*** &       -0.126 &        0.821 &       -1.792*** &       -0.447 &       -0.035 &        1.279** \\ 
 & (       0.215) & (       0.271) & (       0.170) & (       0.350) & (       0.577) & (       0.227) & (       0.308) & (       0.043) & (       0.579) \\ 
\textit{RW p-value} & [ 0.037] & [ 0.037] & [ 0.000] & [ 0.845] & [ 0.591] & [ 0.000] & [ 0.591] & [ 0.845] & [ 0.137] \\ 
  & & & & & & & & & \\ 
Mean (pre-1997) &         2.13 &         3.15 &         0.94 &         4.43 &        16.35 &         2.46 &         3.86 &         0.12 &        13.25 \\ 
Obs &       15,642 &       15,642 &       15,642 &       15,642 &       15,642 &       15,642 &       15,642 &       15,642 &       15,642 \\ 
  & & & & & & & & & \\ 
\underline{\textit{Panel C: Males}}  \\ 
\textit{Intensity 1999 x post} &       -1.074*** &       -1.086*** &       -0.559*** &        0.450 &        0.702 &       -1.532*** &       -0.140 &        0.075 &        2.163*** \\ 
 & (       0.262) & (       0.212) & (       0.159) & (       0.334) & (       0.565) & (       0.261) & (       0.263) & (       0.083) & (       0.641) \\ 
\textit{RW p-value} & [ 0.000] & [ 0.000] & [ 0.003] & [ 0.714] & [ 0.714] & [ 0.000] & [ 0.734] & [ 0.734] & [ 0.004] \\ 
  & & & & & & & & & \\ 
Mean (pre-1997) &         2.85 &         2.04 &         0.89 &         5.28 &        15.38 &         2.69 &         3.21 &         0.42 &        15.73 \\ 
Obs &       15,642 &       15,642 &       15,642 &       15,642 &       15,642 &       15,642 &       15,642 &       15,642 &       15,642 \\ 
  & & & & & & & & & \\ 
\bottomrule
\end{tabular}
    }
    \vspace{-3mm}
    \floatfoot{\deleted{Notes: Progresa intensity data are pooled from 1990 to 2000. All regression models include year fixed effects and municipality fixed effects. Standard errors are clustered at the municipality level. The regression model is weighted by the number of population aged +65. ***p$<0.01$, **p$<0.05$, *p$<0.1$.}\added{Notes: This table reports difference-in-differences estimates of the effect of Progresa enrollment intensity on cause-specific Elder Mortality Rate (EMR) for adults aged 65 and older per 1,000 inhabitants, using the Barham-Rowberry sample (municipalities incorporated into Progresa in 1998 or 1999). Each column corresponds to a specific cause of death. All regressions include municipality and year fixed effects. Standard errors are clustered at the municipality level. ***p$<0.01$, **p$<0.05$, *p$<0.1$.}}
  }
\end{table}
\end{landscape}


\begin{table}[htbp]
\centering
\caption{\deleted{Progresa Impacts on Household Expenditures: Heterogeneity by Elderly Presence (Eligible Households)}\added{Short-Run Effects of Progresa on Household Expenditure by Elderly Presence, Eligible Households}}
\label{tab:expenditures_elderly}
\floatbox[{\captop}]{table}[\FBwidth]
  {}
  {\scalebox{0.9}{
      \begin{tabular}{lcccc}
\toprule
 & (1) & (2) & (3) & (4) \\
 & ln Food & ln Medicine & \% Food & \% Medicine \\
\midrule
Treat $\times$ 1999 (No Elderly)   & 0.121$^{***}$ & 0.039 & $-$0.014 & 0.000 \\
                                    & (0.026) & (0.062) & (0.014) & (0.005) \\[4pt]
Diff: Elderly HH $\times$ 1999     & 0.007 & $-$0.012 & 0.021 & 0.004 \\
                                    & (0.035) & (0.115) & (0.017) & (0.011) \\[4pt]
Observations                        & 22,524 & 22,763 & 22,478 & 22,489 \\
Control Mean (1997)                 & 4.668 & 0.744 & 0.757 & 0.053 \\
\bottomrule
\end{tabular}



    }
    \vspace{-3mm}
    \floatfoot{\deleted{Notes: . ***p$<0.01$, **p$<0.05$, *p$<0.1$.}\added{Notes: This table reports difference-in-differences estimates of the effect of Progresa enrollment intensity on household expenditure outcomes, stratified by the presence of elderly members in the household. The sample is restricted to eligible households. All regressions include municipality and year fixed effects. Standard errors are clustered at the municipality level and shown in parentheses. ***p$<0.01$, **p$<0.05$, *p$<0.1$.}}
  }
\end{table}